\section{Численная проверка}
\label{app:num}

\subsection{Продольное движение}

С безразмерным временем $\theta=\omega_s t$ и $\mu=\omega_\rf/\omega_s$ система~\eqref{eq:set-system} принимает вид
\begin{equation}
\label{eq:num-sys}
\frac{\mathrm{d}\varphi}{\mathrm{d}\theta}=\mu\bigl(\eta_0\delta+\eta_1\delta^2+\Lambda\bigr),
\qquad
\frac{\mathrm{d}\delta}{\mathrm{d}\theta}=-\frac{\sin\varphi}{\mu\,\eta_0}.
\end{equation}
Она интегрировалась точно (с $\sin\varphi$, без замены $\cos\varphi\to1$) методом Дормана--Принса с относительной точностью $10^{-12}$; параметры условные, оба знака $\eta_0$. Среднее $\langle\delta\rangle$ и частота $\Omega$ измерялись за синхротронный период и сравнивались с~\eqref{eq:sol-deq} и~\eqref{eq:sol-omega-sum} при измеренной амплитуде первой гармоники. Во всех случаях $\langle\tau\rangle$ равно нулю с точностью $10^{-17}$.

\begin{table}[htb]
\caption{Полная система~\eqref{eq:num-sys}: счёт и формулы. $\delta_m\approx10^{-2}$.}
\label{tab:num-sol}
\small
\begin{ruledtabular}
\begin{tabular}{rrrrcccc}
$\eta_0$ & $\eta_1$ & $\Lambda$ & $\mu$ & $\langle\delta\rangle$, счёт & $\langle\delta\rangle$, \eqref{eq:sol-deq} & $\Omega/\omega_s-1$, счёт & \eqref{eq:sol-omega-sum} \\
\hline
$-0.5$ & $0$ & $0$ & $10$ & $<10^{-17}$ & $0$ & $-1.5628\times10^{-4}$ & $-1.5624\times10^{-4}$ \\
$0.4$ & $0$ & $0$ & $10$ & $<10^{-17}$ & $0$ & $-1.0001\times10^{-4}$ & $-0.9999\times10^{-4}$ \\
$-0.5$ & $0.5$ & $0$ & $1$ & $5.0002\times10^{-5}$ & $5.0000\times10^{-5}$ & $-4.3232\times10^{-5}$ & $-4.3229\times10^{-5}$ \\
$0.4$ & $-0.8$ & $0$ & $1$ & $1.0002\times10^{-4}$ & $1.0000\times10^{-4}$ & $-1.6771\times10^{-4}$ & $-1.6766\times10^{-4}$ \\
$-0.5$ & $0.5$ & $2\times10^{-5}$ & $1$ & $9.0008\times10^{-5}$ & $9.0000\times10^{-5}$ & $-8.3239\times10^{-5}$ & $-8.3229\times10^{-5}$ \\
$0.4$ & $-0.5$ & $2\times10^{-5}$ & $1$ & $1.2500\times10^{-5}$ & $1.2500\times10^{-5}$ & $-3.6043\times10^{-6}$ & $-3.6037\times10^{-6}$ \\
$-0.5$ & $0.5$ & $0$ & $10$ & $4.9999\times10^{-5}$ & $4.9996\times10^{-5}$ & $-1.9793\times10^{-4}$ & $-1.9790\times10^{-4}$ \\
$-0.5$ & $0$ & $2\times10^{-4}$ & $10$ & $4.0000\times10^{-4}$ & $4.0000\times10^{-4}$ & $-1.5628\times10^{-4}$ & $-1.5624\times10^{-4}$ \\
\end{tabular}
\end{ruledtabular}
\end{table}

Первые две строки --- только поправка фазы; частота совпадает с точной частотой маятника до $10^{-15}$, остаток относительно $-\varphi_m^2/16$ есть член $\varphi_m^4$. Остаток среднего в строках 3, 4, 7 совпадает с членом $\delta_m^4$ из~\eqref{eq:sol-a4}. В строке 6 перекрёстный член ($+6.25\times10^{-5}$) и член $5/12$ ($-6.51\times10^{-5}$) почти сокращаются, и счёт воспроизводит разность.

\subsection{Гладкие электростатические кольца}

Релятивистское движение заряда в поле $E\propto r^{-1}$ и $E\propto R^{-2}$ ($c=m=\rho=1$) интегрировалось при $\gamma_s=1.25$ и $2$. Коэффициенты $\eta_0$, $\eta_1$, извлечённые из периодов круговых орбит, совпадают с таблицей~\ref{tab:es-smooth} с точностью $10^{-8}$ (например, $\eta_1=21.75$ для сферического дефлектора при $\gamma_s=2$). Бетатронный вклад определялся по разности периодов для амплитуд $a$ и $2a$, что исключает краевую ошибку неполного бетатронного периода: $\Lambda/X=0.002$ и $0.001$ для цилиндрического дефлектора (ожидается $0$) и $-0.362$, $-0.753$ для сферического (ожидается $-\beta_s^2=-0.36$, $-0.75$).
